\begin{definition}[Многоцветные числа Рамсея]
\(R(n_1,\ldots,n_k)\) — наименьшее \(t\) такое, что при любой раскраске рёбер \(K_t\)
в \(k\) цветов \(\varphi\colon E(K_t)\to\{1,\ldots,k\}\) найдётся \(i\) и клика \(K_{n_i}\),
все рёбра которой цвета \(i\).
\end{definition}
При \(k=2\) получаем \(R(n_1,n_2)\): граф \((V,\varphi^{-1}(1))\) содержит \(K_{n_1}\),
а граф \((V,\varphi^{-1}(2))\) — его дополнение — \(K_{n_2}\).
\begin{tcolorbox}[colback=cbGrayBack,colframe=cbGray!55]
\textbf{Свойства.} \(R\) симметрично (не меняется при перестановке аргументов) и
монотонно по каждому аргументу. Если некоторое \(n_i=1\), то \(R=1\) (одна вершина —
клика \(K_1\) любого цвета). Если \(n_i=2\), то цвет \(i\) нельзя использовать вовсе, и
\(R(n_1,\ldots,2,\ldots,n_k)=R(n_1,\ldots,\widehat{n_i},\ldots,n_k)\).
\end{tcolorbox}

\begin{theorem}[Гринвуд–Глисон, 1955]
При \(n_1,\ldots,n_k\ge2\)
\[
R(n_1,\ldots,n_k)\le\sum_{i=1}^{k}R(n_1,\ldots,n_i-1,\ldots,n_k)-k+2 .
\]
В лекциях записано более слабое \(R(n_1,\ldots,n_k)\le\sum_iR(n_1,\ldots,n_i-1,\ldots,n_k)\).
\end{theorem}
\begin{proof}
Обозначим \(R_i=R(n_1,\ldots,n_i-1,\ldots,n_k)\) и рассмотрим раскраску \(K_N\),
\(N=\sum R_i-k+2\). Возьмём вершину \(v\); пусть \(A_i\) — множество вершин, соединённых
с \(v\) рёбрами цвета \(i\). Тогда \(\sum|A_i|=N-1=\sum(R_i-1)+1\), и по принципу
Дирихле \(|A_i|\ge R_i\) для некоторого \(i\). По определению \(R_i\) в клике на \(A_i\) есть
\begin{itemize}\tightlist
\item либо клика \(K_{n_i-1}\) цвета \(i\) — с вершиной \(v\) она даёт \(K_{n_i}\) цвета \(i\);
\item либо для некоторого \(j\ne i\) клика \(K_{n_j}\) цвета \(j\).
\end{itemize}
В обоих случаях нужная одноцветная клика найдена.
\end{proof}
\begin{graphillustration}{Шаг доказательства Гринвуда–Глисона при $k=3$: рёбра из $v$ делятся по цветам на $A_1,A_2,A_3$. Хотя бы одно $|A_i|\ge R_i$; внутри него ищем $K_{n_i-1}$ цвета $i$ (с $v$ — $K_{n_i}$) или $K_{n_j}$ другого цвета.}{fig:10-pigeon}
\begin{tikzpicture}[font=\small]
\node[gi selected] (v) at (0,0) {$v$};
\foreach \k/\ang/\sty/\lab in {1/150/gi c1/A_1,2/30/gi c2/A_2,3/270/gi c3/A_3} {
  \node[draw,ellipse,minimum width=2.2cm,minimum height=1.2cm] (E\k) at (\ang:2.6) {};
  \node[gi note] at ($(\ang:2.6)+(1.45,0.45)$) {$\lab$};
  \foreach \d [count=\j] in {-0.6,0,0.6} {\node[gi small] (p\k-\j) at ($(\ang:2.6)+(\d,-0.1)$) {}; \draw[\sty] (v) -- (p\k-\j);}
}
\node[gi note,anchor=west] at (4.2,-1.4) {$\sum|A_i|=N-1=\sum(R_i-1)+1$\\$\Rightarrow\ \exists i:\ |A_i|\ge R_i$};
\end{tikzpicture}
\end{graphillustration}

\begin{cor}
Все числа \(R(n_1,\ldots,n_k)\) конечны.
\end{cor}

\begin{cor}[Мультиномиальная оценка]
\[
R(n_1,\ldots,n_k)\le\binom{n_1+\cdots+n_k-k}{n_1-1,\ \ldots,\ n_k-1}
=\frac{(n_1+\cdots+n_k-k)!}{(n_1-1)!\cdots(n_k-1)!}.
\]
\end{cor}
\begin{proof}
Индукция по \(n_1+\cdots+n_k\). Если какое-то \(n_i=1\), левая часть равна 1, а
мультиномиальный коэффициент не меньше 1. Иначе все \(a_i=n_i-1\ge1\), и
по теореме (в слабой форме) и тождеству для мультиномиальных коэффициентов
\[
\binom{s}{a_1,\ldots,a_k}=\sum_{i=1}^{k}\binom{s-1}{a_1,\ldots,a_i-1,\ldots,a_k},
\qquad s=a_1+\cdots+a_k,
\]
получаем \(R(n_1,\ldots,n_k)\le\sum_iR(\ldots,n_i-1,\ldots)\le\binom{s}{a_1,\ldots,a_k}\).
Тождество — разбор по тому, какой «тип» стоит на первом месте в слове из
\(a_1\) букв первого типа, \ldots, \(a_k\) букв \(k\)-го типа.
\end{proof}
При \(k=2\) это ровно \(\binom{m+n-2}{m-1}\) из билета~\ref{ticket:8}.
\begin{example}
\(R(3,3,3)\le R(2,3,3)+R(3,2,3)+R(3,3,2)-3+2=3\cdot6-1=17\) (здесь
\(R(2,3,3)=R(3,3)=6\)). Это точное значение: \(R(3,3,3)=17\).
Мультиномиальная оценка даёт \(\frac{6!}{2!\,2!\,2!}=90\).
\end{example}

\subsection*{Применение: теорема Шура}
\begin{theorem}[Шур, 1916]
Для любого \(k\ge1\) существует \(n\) такое, что при любой раскраске чисел
\(1,\ldots,n\) в \(k\) цветов найдутся \(x,y,z\) одного цвета с \(x+y=z\).
\end{theorem}
\begin{proof}
Пусть \(N=R(\underbrace{3,\ldots,3}_k)\), \(n=N-1\), и \(\varphi\colon\{1,\ldots,n\}\to\{1,\ldots,k\}\) —
раскраска чисел. Раскрасим рёбра \(K_N\) на вершинах \(1,\ldots,N\) правилом
\(\psi(uv)=\varphi(|u-v|)\). По определению \(N\) есть одноцветный треугольник
\(a<b<c\): \(\varphi(b-a)=\varphi(c-b)=\varphi(c-a)\). Положим \(x=b-a\), \(y=c-b\),
\(z=c-a\): тогда \(x+y=z\), и все три одного цвета (допускается \(x=y\)).
\end{proof}
\begin{graphillustration}{Сведение Шура к числам Рамсея: ребро $uv$ получает цвет числа $|u-v|$. Одноцветный треугольник $a<b<c$ даёт $x=b-a$, $y=c-b$, $z=c-a$ одного цвета, $x+y=z$. Справа — раскраска $1,\ldots,4$ в два цвета без решения (синий контур — цвет 1, двойной серый — цвет 2): оценка $n=R(3,3)-1=5$ точна.}{fig:10-schur}
\begin{tikzpicture}[font=\small]
\node[gi vertex] (a) at (0,0) {$a$};
\node[gi vertex] (b) at (2.4,0) {$b$};
\node[gi vertex] (c) at (1.2,1.8) {$c$};
\draw[gi c1] (a) -- node[below] {$b-a$} (b);
\draw[gi c1] (b) -- node[right] {$c-b$} (c);
\draw[gi c1] (a) -- node[left] {$c-a$} (c);
\begin{scope}[shift={(5.0,0.8)}]
\node[gi vertex,draw=cbBlue,line width=1.4pt] at (0,0) {$1$};
\node[gi selected] at (0.9,0) {$2$};
\node[gi selected] at (1.8,0) {$3$};
\node[gi vertex,draw=cbBlue,line width=1.4pt] at (2.7,0) {$4$};
\node[gi note,anchor=west] at (3.3,0) {$1+1=2$, $1+3=4$, $2+2=4$\\— везде разные цвета};
\end{scope}
\end{tikzpicture}
\end{graphillustration}

